\documentclass[12pt,a4paper]{report}

\usepackage[italian]{babel}
\usepackage[utf8]{inputenc}
\usepackage[T1]{fontenc}
\usepackage{graphicx}
\usepackage[hidelinks]{hyperref}
\usepackage{float}
\usepackage{booktabs}
\usepackage{geometry}
\geometry{margin=2.5cm}

\title{Gestionale Palestra}
\author{Michele Tartabini \and Francesco Napoli \and Michele Raganini}
\date{A.A. 2025/2026}

\begin{document}

\maketitle
\tableofcontents
\newpage

%--------------------------------------------------
\chapter{Descrizione Progetto}
\label{ch:descrizione}

%--------------------------------------------------
\chapter{Glossario dei Termini}
\label{ch:glossario}

\begin{table}[H]
\centering
\begin{tabular}{llll}
\toprule
Termine & Descrizione & Tipo & Sinonimi \\
\midrule
 &  &  &  \\
\bottomrule
\end{tabular}
\end{table}

%--------------------------------------------------
\chapter{Gestione dei Requisiti}
\label{ch:requisiti}

\section{Requisiti Funzionali}

\section{Requisiti Non Funzionali}
RNF1: Implementazione in Python 3 \\
Il sistema dovrà essere implementato in linguaggio Python 3.\\
\\
RNF2: Utilizzo di un database relazionale (SQLite) \\
Il sistema dovrà utilizzare un database relazionale per il salvataggio dei dati.\\
\\
RNF3: Implementazione grafica\\
Il sistema dovrà implementare un'interfaccia grafica in PyQt6.\\
\\
RNF4: Password solo hashate\\ 
Le password non dovranno essere visualizzabili in chiaro ma solo sotto forma di hashcode.\\
\\
RNF5: Interfaccia User-Friendly\\
L'interfaccia dovrà essere completamente in italiano, i messaggi di errore chiari, in modo che un utente abituato a gestionali base non necessiterà di nessun tipo di formazione per usare l'app.\\
\\


\section{Diagrammi dei Casi d'Uso}

\subsection{Deltaglio Casi d'Uso}

\section{Matrice di Mapping}

%--------------------------------------------------
\chapter{Analisi}
\label{ch:analisi}

\section{Diagramma delle Classi di Analisi}

\section{Diagrammi di Sequenza}

\section{Diagrammi di Attivit\`a}

%--------------------------------------------------
\chapter{Progettazione}
\label{ch:progettazione}

\section{Classi di Progettazione}

\section{Diagrammi delle Macchine a Stati}

\section{Diagramma dei Componenti}

\section{Diagramma Entit\`a-Relazioni}

\section{Diagramma di Deployment}

%--------------------------------------------------
\chapter{Implementazione}
\label{ch:implementazione}

\section{Mockup}

\section{Tecnologie Utilizzate}

\section{Unit Test}

\end{document}
